
\section{Einleitung
  \label{kepler:section:Einleitung}}

Johannes Kepler veröffentlichte im frühen 17. Jahrhundert die Gesetze der Planetenbahnen, welche den Grundstein für die heutige Astronomie legten. Kepler stützte seine Erkenntnisse auf die Auswertung der Beobachtungsdaten von Tycho Brahe. Erst viele Jahrzehnte später gelang es Isaac Newton, eine  physikalische Begründung für diese Gesetze zu finden.

Das Haupotthema dieser Arbeit ist die exakte mathematische Herleitung des Dritten Keplerschen Gesetzes. Dieses Gesetz beschreibt den Zusammenhang zwischen der Umlaufzeit eines Planeten und der Grösse seiner Bahn (der grossen Halbachse). Vielfach beschränkt man sich bei der Herleitung oft auf den Spezialfall von Kreisbahnen. Da reale Planetenbahnen jedoch keine Kreise, sondern Ellipsen beschreiben, reicht diese vereinfachte Betrachtung für einen allgemeinen Beweis nicht aus.

Die Herleitung Unterteilen wir in drei Schritte:
 Zunächst werden das erste und zweite Keplersche Gesetz als notwendige physikalische Voraussetzungen formuliert. Anschliessend wird die gängige Herleitung für Kreisbahnen durchgeführt. Im letzten Schritt wird der vereinfachte Ansatz der Kreisbahn durch die allgemeingültigen Ellipsenbahn ersetzt. Die Herleitung machen wir durch die komplexe Integration der Bewegungsgleichungen.